% ============================================================
%  附录 B  人物小传
% ============================================================
\section*{附录 B\quad 人物小传}
\addcontentsline{toc}{section}{附录 B\quad 人物小传}
\markboth{附录 B\ 人物小传}{附录 B\ 人物小传}

\noindent{\small 按书中出现顺序排列。只记与本书主题相关的事实，不作全面评价。}

\vspace{0.6em}

\noindent\textbf{Michael Atiyah（1929--2019）}\quad
英国数学家，1966 年 Fields 奖得主。与 Hirzebruch 创立拓扑 $K$ 理论，与 Singer 证明指标定理；与 Ward、Drinfeld、Hitchin、Manin 构造 ADHM 瞬子解；与 Hitchin 研究磁单极子，合著《磁单极子的几何与动力学》（1988）。他是把物理问题翻译成几何问题的代表人物，也是本书叙事的起点与终点。

\medskip
\noindent\textbf{Edward Witten（1951-- ）}\quad
美国理论物理学家，1990 年 Fields 奖得主（迄今唯一）。弦论、超对称、拓扑场论的奠基者之一；1988 年提出拓扑扭曲与 Donaldson 理论的联系，1994 年与 Seiberg 解出 $\cN=2$ 超杨--米尔斯理论，2006 年与 Kapustin 给出几何 Langlands 的规范理论解释。

\medskip
\noindent\textbf{Peter Goddard、Jean Nuyts、David Olive（GNO）}\quad
1977 年三人在 Nuclear Physics B125 上发表《规范理论与磁荷》，提出磁荷由对偶群的权标号，并指出这个对偶群与 Langlands 对偶群同构。Goddard 后来任剑桥 DAMTP 主任、Newton 研究所所长；Olive 在 CERN 与 Imperial College 工作，是 1970 年代对偶研究的中枢人物。

\medskip
\noindent\textbf{Claus Montonen、David Olive}\quad
1977 年在 Physics Letters B72 上提出“磁单极子作为规范粒子”的猜想：量子化之后，单极子与光子属于同一个规范多重态；两篇论文相隔数月，常被合称 Montonen--Olive 对偶。

\medskip
\noindent\textbf{Hugh Osborn}\quad
剑桥大学应用数学与理论物理系。1979 年计算 $\cN=4$ 超对称规范理论的中心荷，指出只有 $\cN=4$ 的单极子自旋与基本粒子匹配——这是对 Montonen--Olive 猜想的第一次“筛选”。

\medskip
\noindent\textbf{Ashoke Sen（1956-- ）}\quad
印度理论物理学家，Harish-Chandra 研究所教授。1994 年的文章用两单极子模空间上的 $L^2$ 调和形式证明了对偶所预言的束缚态存在，被 Witten 称为“smoking gun”。他是弦论对偶（S 对偶、T 对偶）研究的核心人物之一。

\medskip
\noindent\textbf{Nathan Seiberg（1956-- ）}\quad
美国理论物理学家，普林斯顿高等研究院教授。1993--1994 年在 $\cN=1$ 动力学上取得突破，随后与 Witten 解出 $\cN=2$ 的超杨--米尔斯理论（Seiberg--Witten 理论）。

\medskip
\noindent\textbf{Nigel Hitchin（1946-- ）}\quad
英国数学家，牛津大学教授。与 Atiyah 研究单极子模空间，发现 Atiyah--Hitchin 流形；提出 Hitchin 系统与 Higgs 丛，后者成为几何 Langlands 的几何舞台。

\medskip
\noindent\textbf{Vladimir Drinfeld（1954-- ）}\quad
乌克兰裔美国数学家，1990 年 Fields 奖得主。与 Beilinson 提出几何 Langlands 纲领；在量子群、可积系统、几何 Satake 等方面有奠基性工作。

\medskip
\noindent\textbf{Alexander Beilinson（1957-- ）}\quad
美国数学家，芝加哥大学教授。与 Drinfeld 合著《Quantization of Hitchin's Integrable System and Hecke Eigensheaves》，几何 Langlands 的经典文献。

\medskip
\noindent\textbf{Anton Kapustin（1971-- ）}\quad
俄裔美国理论物理学家，加州理工学院教授。2005 年给出 't Hooft 算子的精确定义；2006 年与 Witten 合写《Electric-Magnetic Duality and the Geometric Langlands Program》，把几何 Langlands 翻译成规范理论的镜像对称。

\medskip
\noindent\textbf{Robert Langlands（1936-- ）}\quad
加拿大数学家，2018 年 Abel 奖得主。1967 年提出 Langlands 纲领，把数论、表示论与自守形式统一在一个猜想框架里。

\medskip
\noindent\textbf{Gerard 't Hooft（1946-- ）、Alexander Polyakov（1945-- ）}\quad
荷兰与俄罗斯理论物理学家，'t Hooft 于 1999 年获 Nobel 物理学奖。1974 年两人各自发现 $SU(2)$ 破缺到 $U(1)$ 时的光滑磁单极子；'t Hooft 1980 年定义 't Hooft 算子，并预言禁闭的对偶迈斯纳机制。

\medskip
\noindent\textbf{E. B. Bogomolny、M. K. Prasad、C. M. Sommerfield}\quad
1975--1976 年给出单极子能量下界与磁荷 $1$ 的解析解（Bogomolny 方程、Prasad--Sommerfield 解）。

\medskip
\noindent\textbf{Peter Kronheimer、Matthias Pauly}\quad
1986 年与 1990 年分别研究带狄拉克型奇点的 Bogomolny 方程，是第 9 章的先声。

\medskip
\noindent\textbf{Simon Donaldson（1957-- ）}\quad
英国数学家，1986 年 Fields 奖得主。1982 年用反自对偶联络构造四维流形不变量；后来与 Kronheimer 合著《四维流形的几何》。

\medskip
\noindent\textbf{M. S. Narasimhan、S. Ramanan}\quad
印度数学家，1960--70 年代研究紧 Riemann 面上的向量丛模空间，引入 Hecke 变换（Hecke 对应），为几何 Langlands 提供了最早的几何构件。

\medskip
\noindent\textbf{Dennis Gaitsgory（1973-- ）}\quad
美国数学家，马克斯·普朗克数学研究所所长。几何 Langlands 纲领的主要推动者；2024 年与十余位合作者宣布证明非分歧情形的几何 Langlands 猜想。
